
The CEDAR-KTAG system delivers kaon identification for the NA62 experiment.
A new \gls{CEDAR} with \gls{h2} as a radiator gas was designed by NA62 group at the University of Birmingham and built at CERN where it was characterised and tested at a test-beam in 2022. 
\gls{ch} was installed and commissioned in the NA62 beam-line for the 2023 run. A number of pressure scans and other tests were completed to ensure that the detector was ready for data taking. 
The detector produced excellent results with 13\,\% greater light yield compared with \gls{cw} with \gls{n2}. \gls{ch} has been in operation since 2023 and has performed flawlessly. The KTAG with \gls{ch} resulted in a kaon identification efficiency of 99.7\,\% and a kaon time resolution of 66\,ps.

The search for the dark photon in $\pi^0$ decays at NA62 was conducted with two studies with different background estimation methods: simulation (MC) and data-driven. The MC study did not produce competitive results in the search region when compared to previous experiments.
Utilising the data-driven method the \gls{UL} of the mixing parameter provides improved coverage across the phase space not covered by any other experiments.
It is planned that this search will continue by analysing data collected in 2025 at NA62 to further improve the result.

The NA62 software package was upgraded to allow for more flexibility. The upgrade was rigorously tested to ensure stability and compatibility with earlier software versions. This allowed for the implementation and testing of new detectors for future experiments such as \gls{HIKE}. The upgrades planned for the KTAG and RICH and their performances within \gls{HIKE} Phase 1 are shown. The 
particle rates for the detectors for \gls{HIKE} Phase 2 were obtained. These results aided the conceptual design of the \gls{HIKE} detectors. The results from the simulations were used in the \gls{HIKE} proposal submitted in 2023 to the CERN SPS committee. Unfortunately the \gls{HIKE} experiment was not approved. However, the upgraded software and detector studies have been used to produce the KOTO-II proposal submitted to the J-PARC committee in 2024. 